\documentclass[10pt,a4paper]{amsart}
\usepackage{amsmath,amssymb,amsthm,graphicx}
\usepackage{verbatim}
\newenvironment{reference}{\comment}{\endcomment}
%\newenvironment{reference}{}{}
\newenvironment{slogan}{\comment}{\endcomment}
\newenvironment{history}{\comment}{\endcomment}

% For commutative diagrams we use Xy-pic
\usepackage[all]{xy}

% We use 2cell for 2-commutative diagrams.
\xyoption{2cell}
% Theorem environments.
%
\theoremstyle{plain}
\newtheorem{theorem}[subsection]{Theorem}
\newtheorem{proposition}[subsection]{Proposition}
\newtheorem{lemma}[subsection]{Lemma}

\theoremstyle{definition}
\newtheorem{definition}[subsection]{Definition}
\newtheorem{example}[subsection]{Example}
\newtheorem{exercise}[subsection]{Exercise}
\newtheorem{situation}[subsection]{Situation}

\theoremstyle{remark}
\newtheorem{remark}[subsection]{Remark}
\newtheorem{remarks}[subsection]{Remarks}

\numberwithin{equation}{subsection}

% Macros
%
\def\lim{\mathop{\mathrm{lim}}\nolimits}
\def\colim{\mathop{\mathrm{colim}}\nolimits}
\def\Spec{\mathop{\mathrm{Spec}}}
\def\Hom{\mathop{\mathrm{Hom}}\nolimits}
\def\Ext{\mathop{\mathrm{Ext}}\nolimits}
\def\SheafHom{\mathop{\mathcal{H}\!\mathit{om}}\nolimits}
\def\SheafExt{\mathop{\mathcal{E}\!\mathit{xt}}\nolimits}
\def\Sch{\mathit{Sch}}
\def\Mor{\mathop{\mathrm{Mor}}\nolimits}
\def\Ob{\mathop{\mathrm{Ob}}\nolimits}
\def\Sh{\mathop{\mathit{Sh}}\nolimits}
\def\NL{\mathop{N\!L}\nolimits}
\def\CH{\mathop{\mathrm{CH}}\nolimits}
\def\proetale{{pro\text{-}\acute{e}tale}}
\def\etale{{\acute{e}tale}}
\def\QCoh{\mathit{QCoh}}
\def\Ker{\mathop{\mathrm{Ker}}}
\def\Im{\mathop{\mathrm{Im}}}
\def\Coker{\mathop{\mathrm{Coker}}}
\def\Coim{\mathop{\mathrm{Coim}}}

% Boxtimes
%
\DeclareMathSymbol{\boxtimes}{\mathbin}{AMSa}{"02}

%
% Macros for moduli stacks/spaces
%
\def\QCohstack{\mathcal{QC}\!\mathit{oh}}
\def\Cohstack{\mathcal{C}\!\mathit{oh}}
\def\Spacesstack{\mathcal{S}\!\mathit{paces}}
\def\Quotfunctor{\mathrm{Quot}}
\def\Hilbfunctor{\mathrm{Hilb}}
\def\Curvesstack{\mathcal{C}\!\mathit{urves}}
\def\Polarizedstack{\mathcal{P}\!\mathit{olarized}}
\def\Complexesstack{\mathcal{C}\!\mathit{omplexes}}
% \Pic is the operator that assigns to X its picard group, usage \Pic(X)
% \Picardstack_{X/B} denotes the Picard stack of X over B
% \Picardfunctor_{X/B} denotes the Picard functor of X over B
\def\Pic{\mathop{\mathrm{Pic}}\nolimits}
\def\Picardstack{\mathcal{P}\!\mathit{ic}}
\def\Picardfunctor{\mathrm{Pic}}
\def\Deformationcategory{\mathcal{D}\!\mathit{ef}}
\usepackage{cleveref}
\crefname{originalsection}{original section}{original sections}
\Crefname{originalsection}{Original section}{Original sections}
\makeatletter
\newlabel{sheaves-section-ringed-spaces}{{0090}{source}{}{sourcepointer-1}{}}
\newlabel{sheaves-section-ringed-spaces@cref}{{[originalsection][0][]0090}{[1][0][]source}{}{}{}}
\makeatother
\crefname{originalsection}{original section}{original sections}
\Crefname{originalsection}{Original section}{Original sections}
\makeatletter
\newlabel{more-algebra-section-injectives-modules}{{01D8}{source}{}{sourcepointer-2}{}}
\newlabel{more-algebra-section-injectives-modules@cref}{{[originalsection][0][]01D8}{[1][0][]source}{}{}{}}
\makeatother
\crefname{originalsection}{original section}{original sections}
\Crefname{originalsection}{Original section}{Original sections}
\makeatletter
\newlabel{sheaves-example-sheaf-product-pointwise}{{006X}{source}{}{sourcepointer-3}{}}
\newlabel{sheaves-example-sheaf-product-pointwise@cref}{{[originalsection][0][]006X}{[1][0][]source}{}{}{}}
\makeatother
\crefname{originalsection}{original section}{original sections}
\Crefname{originalsection}{Original section}{Original sections}
\makeatletter
\newlabel{sheaves-example-sheaf-product-pointwise-algebraic-structure}{{007R}{source}{}{sourcepointer-4}{}}
\newlabel{sheaves-example-sheaf-product-pointwise-algebraic-structure@cref}{{[originalsection][0][]007R}{[1][0][]source}{}{}{}}
\makeatother
\crefname{originallemma}{original lemma}{original lemmas}
\Crefname{originallemma}{Original lemma}{Original lemmas}
\makeatletter
\newlabel{sheaves-lemma-sheaf-subset-stalks}{{0079}{source}{}{sourcepointer-5}{}}
\newlabel{sheaves-lemma-sheaf-subset-stalks@cref}{{[originallemma][0][]0079}{[1][0][]source}{}{}{}}
\makeatother
\crefname{originalsection}{original section}{original sections}
\Crefname{originalsection}{Original section}{Original sections}
\makeatletter
\newlabel{sheaves-section-skyscraper-sheaves}{{0099}{source}{}{sourcepointer-6}{}}
\newlabel{sheaves-section-skyscraper-sheaves@cref}{{[originalsection][0][]0099}{[1][0][]source}{}{}{}}
\makeatother
\crefname{originallemma}{original lemma}{original lemmas}
\Crefname{originallemma}{Original lemma}{Original lemmas}
\makeatletter
\newlabel{sheaves-lemma-stalk-skyscraper-adjoint}{{009C}{source}{}{sourcepointer-7}{}}
\newlabel{sheaves-lemma-stalk-skyscraper-adjoint@cref}{{[originallemma][0][]009C}{[1][0][]source}{}{}{}}
\makeatother
\crefname{originallemma}{original lemma}{original lemmas}
\Crefname{originallemma}{Original lemma}{Original lemmas}
\makeatletter
\newlabel{homology-lemma-adjoint-preserve-injectives}{{015Z}{source}{}{sourcepointer-8}{}}
\newlabel{homology-lemma-adjoint-preserve-injectives@cref}{{[originallemma][0][]015Z}{[1][0][]source}{}{}{}}
\makeatother
\crefname{originallemma}{original lemma}{original lemmas}
\Crefname{originallemma}{Original lemma}{Original lemmas}
\makeatletter
\newlabel{homology-lemma-product-injectives}{{0137}{source}{}{sourcepointer-9}{}}
\newlabel{homology-lemma-product-injectives@cref}{{[originallemma][0][]0137}{[1][0][]source}{}{}{}}
\makeatother
\title{Sheaves of modules on a ringed space have functorial injective embeddings}
\author{The Stacks Project Authors}
\date{stacks-project:injectives.tex:local-snapshot-20261012; GNU FDL 1.2 or later}
\begin{document}
\maketitle
\noindent\textit{Editorial note.} All delivered mathematical text is the author's own, copied from the pinned Stacks Project chapter file sites-cohomology.tex; no statement or proof step was rewritten, reordered or removed.\par\medskip
\noindent\textit{Editorial note.} The proof is the author's complete argument; the construction is displayed and each of its properties is either shown or attributed to a named result carried on this card as a navigation stub.\par\medskip
\noindent\textit{Editorial note.} Every labelled reference occurring inside the counted spans is carried in assembly.external\_references as a non-load-bearing navigation pointer (same-chapter ones at their definition site, cross-chapter ones at the invocation inside the counted span), so no load-bearing mathematics is replaced by a stub. No bibliography citation occurs inside the counted spans.\par\medskip
\noindent\textit{Editorial note.} The same-name check for this card scanned every staged Stacks chapter of this repository -- all sources/stacks-*/pinned-original-source.tex files, the card's own chapter excluded -- and the labels found are recorded in dedup.same\_name\_cross\_chapter on the card.\par\medskip
\noindent\textit{Editorial note.} The preamble is the chapter's own retained minimal preamble (sources/stacks-sites-cohomology/retained-preamble.tex), a verbatim excerpt of the native preamble: the Xy-pic setup and the native theorem declarations with the mathematical macros.\par\medskip
\noindent\textit{Allowed standard prerequisites (editorial).}\begin{itemize}
\item Functorial injective embeddings of modules. More on Algebra, Section more-algebra-section-injectives-modules. Supplies, for every ring and module over it, a functorial injective embedding of the module into an injective module.
\item Pointwise products of sheaves. Sheaves, Examples sheaves-example-sheaf-product-pointwise and sheaves-example-sheaf-product-pointwise-algebraic-structure. Supply that the assignment sending an open set to the product of the chosen injective modules over the stalks is a sheaf of modules.
\item Injectivity detected on stalks. Sheaves, Lemma sheaves-lemma-sheaf-subset-stalks and the description of injective sheaves of modules. Supply that a map of sheaves of modules is injective when it is so on stalks, and that a sheaf constructed stalkwise from injective modules is injective.
\item Ringed spaces. Sheaves, Section sheaves-section-ringed-spaces. Supplies the structure sheaf and the stalks of the ringed space that the construction uses.
\end{itemize}\medskip
\ifdefined\hypertarget\hypertarget{sourcepointer-1}{}\fi
\noindent\textit{Source reference.} Original section 0090: sources/stacks-sheaves/pinned-original-source.tex, source lines 3248--3248. Navigation pointer invoked at proof line 717; the counted span contains the pointer only and the cited result is not reproduced inside it, so the pointer is registered as a non-load-bearing navigation stub in the established form for Stacks material. No bibliography citation is involved.\par\medskip
\ifdefined\hypertarget\hypertarget{sourcepointer-2}{}\fi
\noindent\textit{Source reference.} Original section 01D8: sources/stacks-more-algebra/pinned-original-source.tex, source lines 13617--13617. Navigation pointer invoked at proof line 665; the counted span contains the pointer only and the cited result is not reproduced inside it, so the pointer is registered as a non-load-bearing navigation stub in the established form for Stacks material. No bibliography citation is involved.\par\medskip
\ifdefined\hypertarget\hypertarget{sourcepointer-3}{}\fi
\noindent\textit{Source reference.} Original section 006X: sources/stacks-sheaves/pinned-original-source.tex, source lines 814--814. Navigation pointer invoked at proof line 667; the counted span contains the pointer only and the cited result is not reproduced inside it, so the pointer is registered as a non-load-bearing navigation stub in the established form for Stacks material. No bibliography citation is involved.\par\medskip
\ifdefined\hypertarget\hypertarget{sourcepointer-4}{}\fi
\noindent\textit{Source reference.} Original section 007R: sources/stacks-sheaves/pinned-original-source.tex, source lines 1522--1522. Navigation pointer invoked at proof line 729; the counted span contains the pointer only and the cited result is not reproduced inside it, so the pointer is registered as a non-load-bearing navigation stub in the established form for Stacks material. No bibliography citation is involved.\par\medskip
\ifdefined\hypertarget\hypertarget{sourcepointer-5}{}\fi
\noindent\textit{Source reference.} Original lemma 0079: sources/stacks-sheaves/pinned-original-source.tex, source lines 1163--1163. Navigation pointer invoked at proof line 677; the counted span contains the pointer only and the cited result is not reproduced inside it, so the pointer is registered as a non-load-bearing navigation stub in the established form for Stacks material. No bibliography citation is involved.\par\medskip
\ifdefined\hypertarget\hypertarget{sourcepointer-6}{}\fi
\noindent\textit{Source reference.} Original section 0099: sources/stacks-sheaves/pinned-original-source.tex, source lines 3487--3487. Navigation pointer invoked at proof line 689; the counted span contains the pointer only and the cited result is not reproduced inside it, so the pointer is registered as a non-load-bearing navigation stub in the established form for Stacks material. No bibliography citation is involved.\par\medskip
\ifdefined\hypertarget\hypertarget{sourcepointer-7}{}\fi
\noindent\textit{Source reference.} Original lemma 009C: sources/stacks-sheaves/pinned-original-source.tex, source lines 3549--3549. Navigation pointer invoked at proof line 696; the counted span contains the pointer only and the cited result is not reproduced inside it, so the pointer is registered as a non-load-bearing navigation stub in the established form for Stacks material. No bibliography citation is involved.\par\medskip
\ifdefined\hypertarget\hypertarget{sourcepointer-8}{}\fi
\noindent\textit{Source reference.} Original lemma 015Z: sources/stacks-homology/pinned-original-source.tex, source lines 7487--7487. Navigation pointer invoked at proof line 762; the counted span contains the pointer only and the cited result is not reproduced inside it, so the pointer is registered as a non-load-bearing navigation stub in the established form for Stacks material. No bibliography citation is involved.\par\medskip
\ifdefined\hypertarget\hypertarget{sourcepointer-9}{}\fi
\noindent\textit{Source reference.} Original lemma 0137: sources/stacks-homology/pinned-original-source.tex, source lines 7335--7335. Navigation pointer invoked at proof line 699; the counted span contains the pointer only and the cited result is not reproduced inside it, so the pointer is registered as a non-load-bearing navigation stub in the established form for Stacks material. No bibliography citation is involved.\par\medskip
\begin{lemma}
\label{lemma-sheaves-modules-space}
Let $(X, \mathcal{O}_X)$ be a ringed space, see
Sheaves, Section \ref{sheaves-section-ringed-spaces}.
The category of sheaves of $\mathcal{O}_X$-modules on $X$
has enough injectives. In fact it has functorial injective embeddings.
\end{lemma}

\begin{proof}
For any ring $R$ and any $R$-module $M$ we denote
$j : M \to J_R(M)$ the functorial
injective embedding constructed in
More on Algebra, Section \ref{more-algebra-section-injectives-modules}.
Let $\mathcal{F}$ be a sheaf of $\mathcal{O}_X$-modules on $X$.
By Sheaves, Examples \ref{sheaves-example-sheaf-product-pointwise}
and \ref{sheaves-example-sheaf-product-pointwise-algebraic-structure}
the assignment
$$
\mathcal{I} : U \mapsto
\mathcal{I}(U) = \prod\nolimits_{x\in U} J_{\mathcal{O}_{X, x}}(\mathcal{F}_x)
$$
is an abelian sheaf.
There is a canonical map $\mathcal{F} \to \mathcal{I}$
given by mapping $s \in \mathcal{F}(U)$ to $\prod_{x \in U} j(s_x)$
where $s_x \in \mathcal{F}_x$ denotes the germ of $s$ at $x$.
This map is injective, see
Sheaves, Lemma \ref{sheaves-lemma-sheaf-subset-stalks}
for example.

\medskip\noindent
It remains to prove the following: Given a rule
$x \mapsto I_x$ which assigns to each point $x \in X$ an injective
$\mathcal{O}_{X, x}$-module
the sheaf $\mathcal{I} : U \mapsto \prod_{x \in U} I_x$
is injective. Note that
$$
\mathcal{I} = \prod\nolimits_{x \in X} i_{x, *}I_x
$$
is the product of the skyscraper sheaves $i_{x, *}I_x$ (see
Sheaves, Section \ref{sheaves-section-skyscraper-sheaves} for notation.)
We have
$$
\Hom_{\mathcal{O}_{X, x}}(\mathcal{F}_x, I_x)
=
\Hom_{\mathcal{O}_X}(\mathcal{F}, i_{x, *}I_x).
$$
see Sheaves, Lemma \ref{sheaves-lemma-stalk-skyscraper-adjoint}. Hence it is
clear that each $i_{x, *}I_x$ is an injective $\mathcal{O}_X$-module
(see Homology, Lemma \ref{homology-lemma-adjoint-preserve-injectives} or argue
directly). Hence the injectivity of $\mathcal{I}$ follows from
Homology, Lemma \ref{homology-lemma-product-injectives}.
\end{proof}

\end{document}
